
\documentclass[main.tex]{subfiles}
\begin{document}

	
\chapter{Works of J.-G. Jiang (IV)} \label{ch8}
	
\section{Collectionwise Normal Spaces} \label{sec:8.1}
	

{\definition{\label{def:8.1.1}} \textup{([Bing 1951]\textsuperscript{\cite{Bing}})}. 
$X$ is 
\uwave{$\lambda$-collectionwise normal}
\index{$\lambda$-collectionwise normal| \S\ref{def:8.1.1}\hfill} 
\index{collectionwise normal| \S\ref{def:8.1.1}\hfill} 
if every discrete $\lambda$-family of closed subsets of $X$ has a discrete $($or equivalently, disjoint$)$ open expansion.
}

{\definition{\label{def:8.1.2}}  
$X$ is 
\uwave{$\lambda$-discretely $\varepsilon^*$-expandable}
\index{$\lambda$-discretely $\varepsilon^*$-expandable| \S\ref{def:8.1.2}\hfill} 
\index{discretely $\varepsilon^*$-expandable| \S\ref{def:8.1.2}\hfill} 
if every discrete 
$\lambda$-family of closed subsets of $X$ has a 
$\varepsilon^*$-sequence of open expansions.
}


{\theorem{\label{th:8-1-3}}  \textup{([Jiang 1992]\textsuperscript{\cite{Jiang1992}})}
$X$ is $\lambda$-collectionwise normal iff $X$ is $\lambda$-discretely 
$\varepsilon^*$-expandable.
}

\begin{proof}
$(\Rightarrow)$ is clear.
We prove $(\Leftarrow)$. 
Suppose that $X$ is $\lambda$-discretely $\varepsilon^*$-expandable and 
let $\mathscr{D} = \{D_\alpha: \alpha < \lambda\}$ be a 
discrete closed family in $X$. Let $\{\mathscr{V}_n: n \in \omega\}$ be an
$\varepsilon^*$-sequence of open expansions of $\mathscr{D}$, where each $\mathscr{V}_n = \{V_{n,\alpha}: \alpha < \lambda\}$ with 
$D_\alpha \subset V_{n,\alpha}$ for each 
$n, \alpha$. For $n <\omega$, let 
$H_n = \{x \in X: \vert{}\operatorname{ind}(\mathscr{V}_n)_x\vert{} \le 1\}$. 
Clearly, each $H_n$ is closed in $X$. 
For each $n < \omega$, $\mathscr{V}_n\vert{}_{H_n}$ 
is a point-finite open family in subspace $H_n$. 
It follows that the family 
$\{H_n \bigcap(\bigcup_{\gamma \ne \alpha}V_{n,\gamma}): \alpha < \lambda\}$ 
is an interior-preserving open family in $H_n$. 
For $n \in \omega, \alpha < \lambda$, let
$$B_{n,\alpha} = H_n \setminus \bigcup_{\gamma\ne\alpha} V_{n,\gamma} 
\text{ and } \mathscr{B}_n = \{B_{n,\alpha}: \alpha < \lambda\}.$$

\noindent Then $\mathscr{B} = \bigcup_{n < \omega} \mathscr{B}_n$ is 
a $\sigma$-closure-preserving closed family in $X$. We have

(1) $\bigcup \mathscr{B}^\circ = X$.

\noindent\textsc{Proof}. Let $x \in X$. There exists $m < \omega$ and 
$W \in \operatorname{top}(x)$ such that 
$|\operatorname{ind}(\mathscr{V}_m)_W|\le 1$. 
Let $\alpha_0<\lambda$ be such that 
$\operatorname{ind}(\mathscr{V}_m)_W\subset \{\alpha_0\}$. 
Then $W\subset B_{n,\alpha_0} $, so $x\in B_{m, \alpha_0}^\circ$. Thus (1) holds.

For $n < \omega, \alpha<\lambda$, let

$$U_{n,\alpha} = \operatorname{st}(D_\alpha, \mathscr{B}_n^\circ)
\setminus \bigcup_{i \le n} 
\overline{\bigcup \{B_{i,\gamma}^\circ: \gamma < \alpha
	\text{ and }B_{i,\gamma} \cap D_\alpha = \emptyset\}}$$
and $U_\alpha = \bigcup_{n < \omega} U_{n,\alpha}$. We show that 
$\mathscr{U} = \{U_\alpha: \alpha < \lambda\}$ 
is a disjoint open expansion of $\mathscr{D}$.

(2) $D_\alpha \subset U_\alpha$ for all $\alpha < \lambda$.

\noindent\textsc{Proof}. Let $x \in D_\alpha$. 
By (1), $\exists \, m < \omega, x \in \bigcup \mathscr{B}_m^\circ$, 
$x\in \operatorname{st}(D_\alpha, \mathscr{B}_m^\circ)$.
For each $i \le m$ and $\gamma < \alpha$, 
if $B_{i,\gamma} \cap D_\alpha = \emptyset$, then 
$x \notin B_{i,\gamma}^\circ$. Hence $x \in U_{m,\alpha} \subset U_\alpha$.

(3) If $n < \omega$ and $\gamma\ne \alpha$, then
$$B_{n,\gamma} \cap D_\alpha \subset 
\left(X\setminus \bigcup_{\mu\ne\gamma}V_{n,\mu}\right)
\cap V_{n,\alpha} = \emptyset$$

(4) If $\alpha \neq \beta$, then $U_\alpha \cap U_\beta = \emptyset$.

\noindent\textsc{Proof}.  If not, then $\exists \, x\in U_{\alpha} \cap U_{\beta} $ 
for some 
$\alpha\ne\beta$. Suppose  $x \in U_{m,\alpha} \cap U_{m,\beta}$ and
$m\le n$. Since $x \in U_{m,\alpha}$, there exists 
$\gamma < \lambda$ such that $x \in B_{m,\gamma}^\circ$ and 
$D_\alpha \cap B_{m,\gamma} \neq \emptyset$. 
It follows from (3) that $\gamma = \alpha$, so $x\in B_{m,\alpha}^\circ$. 
On the other hand, since $x\in U_{n, \beta}, m\le n$ and by (3) 
$B_{m,\alpha} \cap D_\beta = \emptyset$, so $x\notin B_{m,\alpha}^\circ$, 
a contradiction. Hence (4) holds.
\end{proof}


{\lemma{} \label{lem:8.1.4}
If a cover $\mathscr{U}$ of $X$ has a $\sigma$-cushioned open refinement, then $\mathscr{U}$ has a cushioned open refinement.
}

\begin{proof}
Suppose that the cover $\mathscr{U} = \{U_\alpha: \alpha \in \Delta\}$ 
has an  open refinement $\mathscr{V} = \bigcup_{n < \omega} \mathscr{V}_n$ 
where  each $\mathscr{V}_n = \{V_{n,\alpha}: \alpha \in \Delta\}$ such that 
for every  $n < \omega$ and $\Gamma \subset \Delta$, 
$\overline{\bigcup_{\alpha \in \Gamma} V_{n,\alpha}} 
\subset \bigcup_{\alpha \in \Gamma} U_\alpha$. 
For $n < \omega$, let $A_n = \bigcup \mathscr{V}_n$, $H_0 = A_0$, and 
$H_n = A_n \setminus \bigcup_{i < n} A_i$ for $n \ge 1$. 
Then $\{H_n: n < \omega\}$ is a 
 locally finite  cover of $X$. For $\alpha \in \Delta$, let 
$B_\alpha = \bigcup_{n < \omega} (H_n \cap V_{n,\alpha})$. 
Then $\mathscr{B} = \{B_\alpha: \alpha \in \Delta\}$ 
is a cushioned refinement of $\mathscr{U}$.
\end{proof}

 {\theorem{\label{th:8-1-5}} 
For any space $X$, the following are equivalent:

\textup{(i)}  $X$ is $\lambda$-collectionwise normal.

\textup{(ii)}  Every $D\omega^2$ $\lambda$-cover of $X$ is a normal cover.

\textup{(iii)}  Every $D\omega^2$ $\lambda$-cover of $X$ has CP closed refinement.
a star refinement.

\textup{(iv)}  Every $D\omega^2$ $\lambda$-cover of $X$ has a semi-open point-star refinement.

\textup{(v)}  Every interior-preserving
$D\omega^2$ $\lambda$-cover of $X$ has a cushioned refinement.
}

\begin{proof}
(i) $\Rightarrow$ (ii). Suppose that $X$ is $\lambda$-collectionwise normal and
 that $\mathscr{U} = \{U_\alpha: \alpha < \lambda\}$ be a $D\omega^2$ 
 $\lambda$-cover of $X$. Let us show that $\mathscr{U}$ is normal.
Let $\mathscr{D} = \bigcup_{m,n < \omega} \mathscr{D}_{m,n}$ be a refinement of
$\mathscr{U}$, where each $\mathscr{D}_{m,n} = \{D_{m,n,\alpha}: \alpha < \lambda\}$ 
is a discrete closed family of subsets of 
$X\setminus \cup\{\cup \mathscr{D}_{p,q}: (p,q) < (m,n)\}$ with 
$D_{m,n,\alpha} \subset U_\alpha$ for each $\alpha < \kappa$. We have 


(i) $\Rightarrow$ (ii) Assume (i) holds and let $\mathscr{U} = \{U_\alpha: \alpha \in \Delta\}$ be an $IPAR(D_0)$ cover of $X$. Let $\{\mathscr{V}_n\}_{n < \omega}$ be a local-star refining sequence for $\mathscr{U}$ by semi-open covers. For $n < \omega$, let

$$A_n = \{x \in X: \mathscr{V}_n \text{ is a local-star refinement of } \mathscr{U} \text{ at } x\}$$

Then $\mathscr{V}_n$ is a local-star refinement of $\mathscr{U}$ at each point of $A_n$. By Lemma 8.2.4, there exists an open family $\mathscr{W}_n$ which is cushioned in $\mathscr{U}$ and covers $A_n$.

Since $\bigcup_{n < \omega} A_n = X$, $\mathscr{W} = \bigcup_{n < \omega} \mathscr{W}_n$ is a $\sigma$-cushioned open refinement of $\mathscr{U}$. By Theorem 8.1.7, this holds.
\end{proof}





\subfile{ch8_2}
\subfile{ch8_3}
%\subfile{ch8_4}

\end{document}
